\chapter*{Preface}
\addcontentsline{toc}{chapter}{Preface}

This book collects, in a single logical order, the twenty papers of the series
\emph{Algorithmic non-commutative class field theory} (September 2026) and the eleven
applied notes that grew out of them. The papers were written in the order in which the
questions arose; the chapters are arranged in the order in which the mathematics is
best learned. Each chapter is a paper, lightly edited: the ``task specification'' framing
of the originals has been left in place where it records a correction, because those
corrections are part of the record, and every chapter keeps its own verification
battery, its own ``Solved / open / impossible'' remark and its own provenance
paragraph. Nothing has been re-proved and nothing has been quietly upgraded.

Three conventions govern the whole book.

\emph{Every number is machine-verified.} Each chapter ends with a table of keyed
checks, each labelled \emph{exact} (integer or symbolic arithmetic), \emph{sampled}
(a statistical test with its $p$-value) or \emph{numeric} (floating point to a stated
tolerance). The scripts that produce the tables, the data they produce, the figures, the
\LaTeX\ sources of every chapter and of this book, and the archived outputs are in the
public repository
\begin{center}
\url{https://github.com/Ruqing1963/anccft}
\end{center}
(\texttt{code/}, \texttt{data/}, \texttt{figures/}, \texttt{papers/},
\texttt{applications/}, \texttt{monograph/}); a frozen copy of each release is archived
on Zenodo under the DOI given in the repository's \texttt{README}. The chapters cite
one another as chapters, never as ``preprints'': nothing in this corpus has been
published elsewhere, and the bibliography lists only the external literature.

\emph{Proof boundaries are stated.} Every result is one of: proved here; quoted from a
named source; verified computationally on named instances; or conjectural. The status
paragraphs at the head of each chapter say which. Where an earlier chapter's claim was
later found to be wrong or imprecise, the correction is recorded in the later chapter and
flagged in the closing chapter of this book.

\emph{The negative results are results.} A large part of the programme consists of
showing that the obvious next step does not exist: no non-negative matrix realizes the
Ihara companion; no commutative Iwasawa algebra governs the congruence tower; no
separable $C^*$-algebra has the pro-module as its $K_0$; no cyclic cocycle computes the
$L$-invariant; no pairwise invariant of a chord diagram computes the arrangement group
of a multi-copy repeat. These are stated as theorems with witnesses, and the closing
chapter lists what remains genuinely open.

The concordance that follows the table of contents maps paper numbers to chapters.
Inside a chapter, ``Theorem 3.7'' means Theorem $n$.3.7 of that chapter $n$; a reference
``[Ch.~1, Thm.~3.7]'' from another chapter points to the same statement.

\bigskip
\noindent Rivi\`ere-Beaudette, September 2026
